\documentclass[12pt]{article}
\usepackage[margin=0.75in]{geometry}

\newcommand{\duedate}{02/10/2026}
\newcommand{\assignment}{MEMORYLLM}

\newcommand{\name}{Aksel Kretsinger-Walters, adk2164}
\newcommand{\email}{adk2164@columbia.edu}

\makeatletter
\def\input@path{{../}{../../}{../../../}}
\makeatother
\input{pset_template_CLMM.tex} %% DO NOT CHANGE THIS LINE

\begin{document}
\psetheader %% DO NOT CHANGE THIS LINE

\section*{MEMORYLLM}

\paragraph{Motivation.}
Large language models encode knowledge implicitly in their parameters and operate within fixed context
windows, limiting their ability to maintain persistent state or accumulate long-term memory across interactions.

\paragraph{Core Idea.}
MEMORYLLM introduces explicit, learnable memory slots that persist across inputs, allowing the model to
maintain and update a structured internal state separate from transient token representations.

\paragraph{Memory Mechanism.}
A fixed set of memory vectors is jointly processed with input tokens via attention, and read/write operations
update these memory slots over time, enabling continuous state evolution beyond a single context window.

\paragraph{Separation of Roles.}
Token embeddings handle local processing and immediate context, while memory vectors store longer-term
information, preventing interference between short-term computation and persistent knowledge.

\paragraph{Persistent State.}
Unlike segment-level recurrence mechanisms, MEMORYLLM maintains memory across sessions or extended
interactions, supporting multi-step reasoning, tool use, and stateful behavior.

\paragraph{Relation to Continual Learning.}
The architecture resembles continual learning systems in that it accumulates information incrementally in
explicit memory rather than solely modifying model parameters, thereby reducing destructive interference.

\paragraph{Comparison to Recurrent Memory Transformer.}
Where RMT focuses on compressing long sequences through segment recurrence, MEMORYLLM emphasizes long-lived,
persistent memory suitable for agent-like applications and interactive settings.

\paragraph{Key Takeaway.}
MEMORYLLM reframes large language models as stateful systems with explicit memory substrates, separating
transient attention-based computation from persistent storage and enabling scalable, long-term reasoning.

\end{document}
